%==============================================================================
% CHAPTER 11: Biological Charge Current and Flux
%==============================================================================

\chapter{Biological Charge Current and Flux}
\label{chap:charge_current}

\epigraph{%
  The membrane is not a wall.\\
  It is a controlled boundary,\\
  and life is the art of controlling it.%
}{after Peter Mitchell, Nobel Lecture (1978)}

\begin{chapterbox}
Charge current in biological systems is not an accident of ionic chemistry.
It is the direct consequence of partition boundary dynamics: when the boundary
$\partial\Part$ moves or deforms, charge flows. Ion channels are controlled
partition apertures with conductance $g = (q^2/h)\,b^{-\Depth_{\mathrm{channel}}}$
--- a function of topology, not molecular detail. The Goldman--Hodgkin--Katz
(GHK) voltage equation emerges from partition theory with permeability
$P_X = b^{-\Depth_X}$; the resting potential $\approx -70$\,mV corresponds
to a specific partition topology of the resting membrane. The Nernst potential
is a partition depth difference $\Delta\Depth$ expressed in voltage units.
The action potential is a partition phase transition: at threshold,
the topology of voltage-gated \ce{Na+} channels undergoes a discontinuous
change, creating an avalanche of charge flow. ATP synthase is a partition
converter that maps a proton-gradient partition depth $\Delta\Depth_{\ce{H+}}$
into the chemical partition depth stored in ATP, with stoichiometry
$n_{\ce{H+}} = 3$ protons per ATP.
\end{chapterbox}

%------------------------------------------------------------------------------
\section{Charge Current from Partition Boundary Motion}
\label{sec:charge_from_boundary}
%------------------------------------------------------------------------------

Chapter~\ref{chap:charge_emergence} established the Charge Emergence Theorem:
electric charge $Q$ is not a primitive quantity but emerges from the topology
of the partition boundary $\partial\Part$. When a closed surface $\partial\Part$
encloses a region of partition depth $\Depth$, the enclosed charge satisfies
\begin{equation}
  Q = \varepsilon_0 \oint_{\partial\Part} \mathbf{E} \cdot d\mathbf{A},
  \label{eq:gauss_partition}
\end{equation}
which is Gauss's law reread as a partition identity: charge is the flux of
the electric field through the partition boundary.

\begin{theorem}[Partition Boundary Current]
\label{thm:partition_boundary_current}
The charge current density $J_{\mathrm{charge}}$ through a biological
membrane equals the rate of change of the electric flux through the
partition boundary:
\begin{equation}
  J_{\mathrm{charge}} = \varepsilon_0\,\frac{\partial}{\partial t}
    \oint_{\partial\Part} \mathbf{E} \cdot d\mathbf{A}.
  \label{eq:partition_current}
\end{equation}
Charge flows if and only if the partition boundary moves, deforms, or the
charge-carrier distribution on it changes.
\end{theorem}

\begin{proof}
Differentiate Eq.~\eqref{eq:gauss_partition} with respect to time:
\begin{equation}
  \frac{dQ}{dt}
    = \varepsilon_0\,\frac{d}{dt}\oint_{\partial\Part} \mathbf{E} \cdot d\mathbf{A}.
\end{equation}
By definition of current, $dQ/dt = I$, and $J_{\mathrm{charge}} = I/A$
through the membrane surface area $A$. The time derivative of the surface
integral changes only when $\partial\Part$ moves (geometric deformation),
or when $\mathbf{E}$ changes on a fixed $\partial\Part$ due to ion
redistribution. In either case the change is driven by partition boundary
dynamics. If $\partial\Part$ is static and the charge distribution is
stationary, $dQ/dt = 0$ and no current flows. \qed
\end{proof}

The physical content is immediate: biological current is \emph{partition
boundary current}. Ion channels, electrogenic pumps, and co-transporters
do not carry charge in any primitive sense; they open or deform partition
boundaries, and charge flows as a purely geometric consequence.

%------------------------------------------------------------------------------
\section{Ion Channels as Controlled Partition Apertures}
\label{sec:ion_channels}
%------------------------------------------------------------------------------

An ion channel is a protein-lined pore spanning the plasma membrane.
In the \emph{closed} state the pore is occluded: the partition depth of the
transmembrane pathway is effectively infinite, $\Depth_{\mathrm{channel}} \to \infty$,
and the partition transition rate (Chapter~\ref{chap:partition_depth},
Theorem on transition kinetics) vanishes. In the \emph{open} state the pore
presents a well-defined partition pathway of depth $\Depth_{\mathrm{channel}}$,
set by the geometry of the selectivity filter and gate.

\begin{definition}[Partition Transmission Coefficient]
\label{def:transmission_coefficient}
The \emph{partition transmission coefficient} of an ion channel is
\begin{equation}
  T_{\mathrm{channel}} = b^{-\Depth_{\mathrm{channel}}},
  \label{eq:transmission_coefficient}
\end{equation}
where $b$ is the partition base (Chapter~\ref{chap:partition_depth}) and
$\Depth_{\mathrm{channel}} \geq 0$ is the partition depth of the open-channel
pathway. A perfectly conducting (zero-depth) channel has $T_{\mathrm{channel}} = 1$.
\end{definition}

The single-channel current follows from treating the open channel as a partition
aperture through which charge carriers transit:
\begin{equation}
  \boxed{I_{\mathrm{channel}}
    = \frac{q^2}{h}\,T_{\mathrm{channel}}\,(V - V_{\mathrm{rev}}),}
  \label{eq:channel_current}
\end{equation}
where $q$ is the ion charge, $h$ is Planck's constant, $V$ is the membrane
voltage, and $V_{\mathrm{rev}}$ is the reversal (Nernst) potential for the
permeant ion. The prefactor $q^2/h$ is the quantum of conductance, appearing
here because the minimum partition-transition time is
$\tau_p = \hbar/(\kB T)$ (Chapter~\ref{chap:bounded_phase_space}), which
yields the fundamental conductance scale
$q^2/h \approx 3.87 \times 10^{-5}$\,S.

\begin{corollary}[Conductance from Partition Topology]
\label{cor:conductance}
The single-channel conductance is
\begin{equation}
  g_{\mathrm{channel}} = \frac{q^2}{h}\,b^{-\Depth_{\mathrm{channel}}}.
  \label{eq:channel_conductance}
\end{equation}
Conductance is a function of partition topology $\Depth_{\mathrm{channel}}$
only; it does not depend on the detailed molecular structure of the channel
protein.
\end{corollary}

\begin{remark}
Observed single-channel conductances span from $\sim 1$\,pS (MscS,
low-conductance mechanosensitive channel) to $\sim 300$\,pS (BK large-conductance
channels). Using Eq.~\eqref{eq:channel_conductance} with base $b = e$:
\begin{align}
  g = 1\ \mathrm{pS}   &\Rightarrow \Depth_{\mathrm{channel}} = \ln(3.87 \times 10^{-5}/10^{-12}) \approx 17.4,\\
  g = 300\ \mathrm{pS} &\Rightarrow \Depth_{\mathrm{channel}} = \ln(3.87 \times 10^{-5}/3\times10^{-10}) \approx 11.7.
\end{align}
The 300-fold conductance range corresponds to a partition depth difference
$\Delta\Depth \approx 5.7$, reflecting roughly six additional partition barriers
in the low-conductance channel --- consistent with the added selectivity-filter
complexity of mechanosensitive channels relative to BK channels.
\end{remark}

\subsection{Gating as Partition Topology Switching}
\label{ssec:gating}

Gating (the opening and closing of an ion channel) is a
\emph{partition topology switch}: $\Depth_{\mathrm{channel}}$ jumps
discontinuously between a large value in the closed state and a small value
$\Depth_{\mathrm{open}}$ in the open state. The gating transition rate is
\begin{equation}
  k_{\mathrm{open}}(V) = \tau_p^{-1}\,b^{-\Delta\Depth_{\mathrm{gate}}(V)},
  \label{eq:gating_rate}
\end{equation}
where $\Delta\Depth_{\mathrm{gate}}(V)$ is the partition depth of the gating
transition, which depends on membrane voltage through the coupling of the
voltage-sensor domain to the gate topology.

For voltage-gated \ce{Na+} channels, the observed activation curve is
Boltzmann-distributed:
\begin{equation}
  P_{\mathrm{open}}(V) = \frac{1}{1 + e^{-(V - V_{1/2})/k_s}},
  \label{eq:boltzmann_gate}
\end{equation}
with $V_{1/2} \approx -40$\,mV and slope $k_s \approx 7$\,mV.
In partition terms:
$\Delta\Depth_{\mathrm{gate}}(V) = \Delta\Depth_0 - qV/(k_s \ln b)$.
The voltage sensitivity of gating is a linear modulation of partition depth
by the transmembrane electric field.

%------------------------------------------------------------------------------
\section{The Goldman--Hodgkin--Katz Equation from Partition Theory}
\label{sec:ghk}
%------------------------------------------------------------------------------

The Goldman--Hodgkin--Katz (GHK) voltage equation is the canonical expression
for the resting membrane potential \citep{Goldman1943,Hodgkin1952}. Its
standard derivation assumes a constant electric field across the membrane and
treats ionic permeability as an empirical fitting parameter. Partition theory
derives both the equation and the permeability from first principles.

\begin{theorem}[GHK Equation from Partition Theory]
\label{thm:ghk}
For a membrane permeable to \ce{K+}, \ce{Na+}, and \ce{Cl-}, the zero-current
membrane voltage is
\begin{equation}
  V_m = \frac{RT}{F}\ln\!\left(
    \frac{P_K[\ce{K+}]_o + P_{Na}[\ce{Na+}]_o + P_{Cl}[\ce{Cl-}]_i}
         {P_K[\ce{K+}]_i + P_{Na}[\ce{Na+}]_i + P_{Cl}[\ce{Cl-}]_o}
  \right),
  \label{eq:ghk}
\end{equation}
where the ionic permeability $P_X$ equals the partition transmission
coefficient for species $X$:
\begin{equation}
  P_X = b^{-\Depth_X},
  \label{eq:permeability_partition}
\end{equation}
and $\Depth_X$ is the partition depth of the transmembrane pathway for ion
species $X$ in the resting configuration.
\end{theorem}

\begin{proof}
The Nernst--Planck equation for the flux of species $X$ (charge $z_X$) across
a membrane of thickness $d$ under a constant electric field $E = V/d$ is
\begin{equation}
  J_X = -D_X\!\left(\frac{d[X]}{dx} + \frac{z_X F}{RT}[X]\,E\right),
  \label{eq:nernst_planck}
\end{equation}
where $D_X$ is the diffusion coefficient within the membrane. Solving
Eq.~\eqref{eq:nernst_planck} with boundary conditions $[X](0) = \beta_X[X]_o$
and $[X](d) = \beta_X[X]_i$ (partition coefficients $\beta_X$) gives the
GHK flux equation:
\begin{equation}
  J_X = P_X\,\frac{z_X^2 F^2 V}{RT}
    \cdot\frac{[X]_i - [X]_o\,e^{-z_X FV/RT}}{1 - e^{-z_X FV/RT}},
  \label{eq:ghk_flux}
\end{equation}
with $P_X = D_X \beta_X / d$. In partition theory, the membrane diffusion
coefficient $D_X$ inside the channel pore is set by the partition transition
rate: $D_X \propto b^{-\Depth_X}$ (each partition boundary crossed reduces
the effective diffusivity by a factor $b^{-1}$). The partition coefficient
$\beta_X$ is unity for a fully permeable boundary and $b^{-\Delta\Depth_X^{\partial}}$
for a boundary of depth $\Delta\Depth_X^{\partial}$. Absorbing all factors
into $P_X = b^{-\Depth_X}$ (where $\Depth_X$ is the total transmembrane
partition depth for species $X$) gives Eq.~\eqref{eq:permeability_partition}.
Setting $\sum_X z_X J_X = 0$ (zero net current at resting potential) and
solving for $V$ yields Eq.~\eqref{eq:ghk}. \qed
\end{proof}

\subsection{Deriving the Resting Potential from Partition Topology}
\label{ssec:resting_potential}

The resting membrane potential of a mammalian neuron, $V_m \approx -70$\,mV,
arises from differential partition depth of the three main ion species.
The resting membrane has partition depths satisfying
$\Depth_K < \Depth_{Cl} < \Depth_{Na}$,
making it far more permeable to \ce{K+} than to \ce{Na+}.
The permeability ratio is
\begin{equation}
  \frac{P_K}{P_{Na}}
    = \frac{b^{-\Depth_K}}{b^{-\Depth_{Na}}}
    = b^{\Depth_{Na} - \Depth_K} \approx 100\ \text{(resting neuron)}.
  \label{eq:permeability_ratio}
\end{equation}
With $b = e$, this gives $\Depth_{Na} - \Depth_K = \ln(100) \approx 4.6$:
the resting membrane imposes approximately 4--5 additional partition barriers
on \ce{Na+} relative to \ce{K+}.

\begin{example}[Resting Potential Calculation]
\label{ex:resting_potential}
Representative mammalian neuronal concentrations:
$[\ce{K+}]_i = 140$\,mM, $[\ce{K+}]_o = 5$\,mM;
$[\ce{Na+}]_i = 15$\,mM, $[\ce{Na+}]_o = 145$\,mM;
$[\ce{Cl-}]_i = 10$\,mM, $[\ce{Cl-}]_o = 110$\,mM.
With $P_K : P_{Na} : P_{Cl} = 1 : 0.04 : 0.45$:
\begin{align}
  V_m &= 25.7\,\mathrm{mV} \cdot \ln\!\left(
    \frac{1\cdot5 + 0.04\cdot145 + 0.45\cdot10}
         {1\cdot140 + 0.04\cdot15 + 0.45\cdot110}
  \right) \notag \\
  &= 25.7\,\mathrm{mV} \cdot \ln\!\left(\frac{15.3}{190.1}\right)
   = 25.7\,\mathrm{mV} \cdot (-2.52)
   \approx -65\ \mathrm{mV}.
\end{align}
The residual $\approx -5$\,mV discrepancy from the physiological value of
$-70$\,mV is accounted for by the electrogenic contribution of the
\ce{Na+}/\ce{K+}-ATPase (Section~\ref{sec:naka_atpase}), which adds
$-3$ to $-5$\,mV through active partition depth maintenance.
\end{example}

Table~\ref{tab:ghk_partition} shows the correspondence between the GHK
permeability parameters and partition depths for three standard ion species.

\begin{table}[htbp]
\centering
\caption{Partition depths $\Depth_X$ (base-$e$ units) corresponding to the
  GHK permeability ratios in the mammalian resting neuron.
  Partition depths are computed via $P_X = e^{-\Depth_X}$,
  normalised to $\Depth_K = 0$ for clarity.}
\label{tab:ghk_partition}
\begin{tabular}{lcccc}
\toprule
Ion species $X$ & $P_X / P_K$ & $\Depth_X - \Depth_K$ & $V_{\mathrm{Nernst}}(X)$ (mV) \\
\midrule
\ce{K+}  & 1.00  & 0.0  & $-89$ \\
\ce{Cl-} & 0.45  & 0.80 & $-63$ \\
\ce{Na+} & 0.04  & 3.22 & $+61$ \\
\bottomrule
\end{tabular}
\end{table}

%------------------------------------------------------------------------------
\section{The Nernst Potential as Partition Depth Difference}
\label{sec:nernst}
%------------------------------------------------------------------------------

For a single ion species $X$ with valence $z$, the Nernst potential is
\begin{equation}
  V_{\mathrm{Nernst}}(X) = \frac{RT}{zF}\ln\frac{[X]_o}{[X]_i}.
  \label{eq:nernst}
\end{equation}
Partition theory gives this equation a geometric meaning.

\begin{theorem}[Nernst Potential as Partition Depth Difference]
\label{thm:nernst_partition}
The Nernst potential for ion species $X$ with valence $z$ is
\begin{equation}
  V_{\mathrm{Nernst}}(X) = \frac{\kB T}{zq}\,\Delta\Depth_X\,\ln b,
  \label{eq:nernst_partition}
\end{equation}
where $\Delta\Depth_X = \Depth_X^{\mathrm{in}} - \Depth_X^{\mathrm{out}}$ is
the partition depth difference for species $X$ across the membrane.
\end{theorem}

\begin{proof}
The concentration ratio encodes a partition depth difference. In partition
theory, the chemical potential of species $X$ in a compartment of concentration
$[X]$ is $\mu_X = \mu_X^0 + \kB T\ln[X]$, and the accessible partition depth
is inversely related to concentration: higher concentration corresponds to
lower partition depth (more accessible states). Specifically,
\begin{equation}
  \frac{[X]_o}{[X]_i} = b^{\Depth_X^{\mathrm{in}} - \Depth_X^{\mathrm{out}}}
    = b^{\Delta\Depth_X}.
  \label{eq:concentration_partition}
\end{equation}
Substituting into Eq.~\eqref{eq:nernst}:
\begin{equation}
  V_{\mathrm{Nernst}}(X)
    = \frac{\kB T}{zq}\ln(b^{\Delta\Depth_X})
    = \frac{\kB T}{zq}\,\Delta\Depth_X\,\ln b.
\end{equation}
The voltage is the partition depth difference across the membrane, rescaled
by the thermal voltage $\kB T/(zq) \approx 25.7$\,mV at $T = 310$\,K for
monovalent ions. \qed
\end{proof}

Table~\ref{tab:nernst_potentials} collects the Nernst potentials and
corresponding partition depth differences for the principal physiological ions.

\begin{table}[htbp]
\centering
\caption{Nernst potentials and partition depth differences for physiological
  ions at $T = 310$\,K. Concentrations are representative mammalian neuronal
  values. $\Delta\Depth_X$ is in base-$e$ units.}
\label{tab:nernst_potentials}
\begin{tabular}{lcccr}
\toprule
Ion & $[X]_i$ (mM) & $[X]_o$ (mM) & $V_{\mathrm{Nernst}}$ (mV) & $\Delta\Depth_X$ \\
\midrule
\ce{K+}       & 140    & 5    & $-89$  & $-3.33$ \\
\ce{Na+}      & 15     & 145  & $+61$  & $+2.27$ \\
\ce{Ca^{2+}}  & 0.0001 & 2    & $+122$ & $+9.49$ \\
\ce{Cl-}      & 10     & 110  & $-63$  & $+2.40$ \\
\bottomrule
\end{tabular}
\end{table}

The extreme Nernst potential of \ce{Ca^{2+}} ($+122$\,mV) reflects the very
large partition depth difference maintained by the cell for this ion:
$|\Delta\Depth_{\ce{Ca^{2+}}}| \approx 9.5$ at rest. This is not accidental.
Calcium serves as a second messenger precisely because the cell maintains an
enormous partition depth reservoir for it; opening a \ce{Ca^{2+}} channel
releases a large partition depth gradient, triggering downstream signalling
cascades with high signal-to-noise ratio.

%------------------------------------------------------------------------------
\section{Partition Flux Conservation}
\label{sec:flux_conservation}
%------------------------------------------------------------------------------

\begin{theorem}[Partition Flux Conservation]
\label{thm:flux_conservation}
In a closed biological system, the total partition flux satisfies the
continuity equation
\begin{equation}
  \oint_{\partial V} \mathbf{J} \cdot d\mathbf{A}
    = -\frac{\partial}{\partial t}\int_V \rho_{\Depth}\,dV,
  \label{eq:partition_flux_conservation}
\end{equation}
where $\rho_{\Depth}$ is the partition depth density (partition depth per
unit volume) and $\mathbf{J}$ is the partition current density (flux of
charge carriers driven by partition depth gradients).
This is the biological analogue of charge conservation.
\end{theorem}

\begin{proof}
From the Charge Emergence Theorem, $Q = \varepsilon_0\oint\mathbf{E}\cdot d\mathbf{A}$.
From Theorem~\ref{thm:nernst_partition}, the electric field generated by a
partition depth gradient is
$\mathbf{E} = -(\kB T \ln b / q)\,\nabla\Depth$.
Combining:
\begin{equation}
  Q = -\varepsilon_0\frac{\kB T \ln b}{q}
    \oint_{\partial V}\nabla\Depth \cdot d\mathbf{A}
    = -\varepsilon_0\frac{\kB T \ln b}{q}
    \int_V\nabla^2\Depth\,dV.
\end{equation}
Charge conservation ($dQ/dt = -\oint\mathbf{J}\cdot d\mathbf{A}$) then implies
\begin{equation}
  \oint_{\partial V}\mathbf{J}\cdot d\mathbf{A}
    = \frac{\partial}{\partial t}\varepsilon_0\frac{\kB T \ln b}{q}
    \int_V\nabla^2\Depth\,dV
    = -\frac{\partial}{\partial t}\int_V\rho_{\Depth}\,dV,
\end{equation}
where we identify $\rho_{\Depth} = -\varepsilon_0(\kB T\ln b/q)\,\nabla^2\Depth$
(Poisson's equation in partition variables). \qed
\end{proof}

\begin{remark}
Equation~\eqref{eq:partition_flux_conservation} is the statement that partition
depth inside any closed volume changes only when partition current flows through
its boundary. A cell that seals all partition boundaries from the environment
conserves its total partition depth --- it is energetically isolated and
therefore dead. Life requires open partition systems that continuously exchange
partition depth with the environment.
\end{remark}

%------------------------------------------------------------------------------
\section{Membrane Potential Dynamics}
\label{sec:membrane_dynamics}
%------------------------------------------------------------------------------

The time evolution of the membrane potential at the single-compartment level is
governed by the balance between capacitive charging and ionic currents:
\begin{equation}
  C_m\frac{dV}{dt}
    = -\sum_X I_X
    = -\sum_X g_X(V - V_X),
  \label{eq:cable_equation}
\end{equation}
where $C_m \approx 1\,\mu\mathrm{F/cm}^2$ is the membrane capacitance,
$g_X = (q^2/h)\,b^{-\Depth_X}$ is the conductance of channel type $X$
(Corollary~\ref{cor:conductance}), and $V_X = V_{\mathrm{Nernst}}(X)$.

\begin{proposition}[Membrane Seeks Minimum Partition Depth]
\label{prop:membrane_minimization}
Equation~\eqref{eq:cable_equation} is equivalent to
\begin{equation}
  C_m\frac{dV}{dt} = -\frac{\partial\Depth_{\mathrm{total}}}{\partial t},
  \label{eq:membrane_depth_flow}
\end{equation}
where $\Depth_{\mathrm{total}} = \sum_X\Depth_X\cdot n_X$ is the total
partition depth summed over all ion species, weighted by the number of open
channels $n_X$ of each type.
\end{proposition}

\begin{proof}
From Corollary~\ref{cor:conductance}, $g_X = (q^2/h)\,b^{-\Depth_X}$.
From Theorem~\ref{thm:nernst_partition}, $V_X = (\kB T/q)\,\Delta\Depth_X\ln b$.
The restoring force on $V$ toward $V_X$ is therefore a restoring force on the
membrane's partition depth toward $\Delta\Depth_X$. Summing,
$\partial\Depth_{\mathrm{total}}/\partial t = -C_m\,dV/dt$, which is
Eq.~\eqref{eq:membrane_depth_flow}. \qed
\end{proof}

The membrane thus acts as a \emph{partition depth minimiser}: it adjusts its
voltage to reduce the total partition depth gradient across the lipid bilayer.
At steady state ($dV/dt = 0$), the resting potential is the conductance-weighted
mean of the individual Nernst potentials:
\begin{equation}
  V_m^{\mathrm{ss}}
    = \frac{\sum_X g_X\,V_X}{\sum_X g_X}
    = \frac{\sum_X b^{-\Depth_X}\,V_X}{\sum_X b^{-\Depth_X}}.
  \label{eq:steady_state_potential}
\end{equation}

%------------------------------------------------------------------------------
\section{Action Potential as a Partition Phase Transition}
\label{sec:action_potential}
%------------------------------------------------------------------------------

The action potential is a $\sim$100\,mV excursion lasting $\sim$1\,ms,
self-propagating along axons. Hodgkin and Huxley \citeyearpar{Hodgkin1952b}
demonstrated that it arises from the sequential voltage-dependent activation
and inactivation of \ce{Na+} and \ce{K+} channels. Partition theory identifies
this sequence as a \emph{partition phase transition}: a cooperative, discontinuous
change in the topology of the voltage-gated channel ensemble at a critical
membrane voltage.

\begin{definition}[Critical Partition Depth]
\label{def:critical_depth}
The critical partition depth $\Depth_{\mathrm{crit}}$ of voltage-gated
\ce{Na+} channels is the depth at which spontaneous channel opening is
thermodynamically inevitable:
\begin{equation}
  \Depth_{\mathrm{Na}}(V_{\mathrm{th}}) = \Depth_{\mathrm{crit}},
  \label{eq:critical_depth_def}
\end{equation}
where $V_{\mathrm{th}} \approx -55$\,mV is the action potential threshold.
At $V = V_{\mathrm{th}}$, the free energy cost of the gating transition equals
$\kB T$, and thermal fluctuations drive the transition with probability one.
\end{definition}

\begin{theorem}[Action Potential as Partition Phase Transition]
\label{thm:action_potential_transition}
The action potential is the following sequence of partition topology changes:
\begin{enumerate}[label=(\roman*)]
  \item \textbf{Threshold crossing.} $V$ reaches $V_{\mathrm{th}}$;
        $\Depth_{\mathrm{Na}}$ reaches $\Depth_{\mathrm{crit}}$. The
        voltage-gated \ce{Na+} channel ensemble undergoes a cooperative
        topology switch from high-depth (closed) to low-depth (open).
  \item \textbf{Avalanche.} $\Depth_{\mathrm{Na}}$ decreases discontinuously.
        \ce{Na+} current $I_{\mathrm{Na}} = g_{\mathrm{Na}}(V - V_{\mathrm{Na}})$
        increases sharply ($V_{\mathrm{Na}} \approx +61$\,mV), driving $V$
        toward $V_{\mathrm{Na}}$.
  \item \textbf{Inactivation.} \ce{Na+} channels enter the inactivated state
        (a third partition topology, $\Depth_{\mathrm{Na}}^{\mathrm{inact}} \gg \Depth_{\mathrm{crit}}$).
        Simultaneously, voltage-gated \ce{K+} channels open ($\Depth_K$ decreases),
        driving $V$ back toward $V_K \approx -89$\,mV.
  \item \textbf{Refractory period.} Both channel types maintain high partition
        depth for their respective pathways, preventing immediate re-excitation
        and establishing the absolute and relative refractory periods.
\end{enumerate}
The threshold behaviour is a genuine phase transition: below $V_{\mathrm{th}}$
the system is monostable (resting state); above $V_{\mathrm{th}}$ it is
transiently monostable toward the action potential peak.
\end{theorem}

The Hodgkin--Huxley model \citep{Hodgkin1952b} is the quantitative realization
of this theorem. Its gating variables $m$, $h$, $n$ are partition occupation
fractions: $m(V) = 1 - b^{-\Delta\Depth_m(V)}$, and similarly for $h$ and $n$.
The cooperativity of \ce{Na+} channel opening (the $m^3 h$ dependence of sodium
conductance) reflects three voltage-sensor domains that must independently
undergo the partition topology switch; each contributes an independent factor
$b^{-\Delta\Depth_{\mathrm{sensor}}(V)}$ to the total transmission coefficient.

%------------------------------------------------------------------------------
\section{ATP Synthase as Partition Converter}
\label{sec:atp_synthase}
%------------------------------------------------------------------------------

ATP synthase (Complex~V of the respiratory chain) converts the
proton-motive force --- a partition depth gradient across the inner
mitochondrial membrane --- into the chemical partition depth stored in ATP.
It is the canonical example of a \emph{partition converter}: a molecular
machine that maps one form of partition depth into another without dissipating
it as heat.

\begin{theorem}[ATP Synthesis as Partition Depth Conversion]
\label{thm:atp_synthesis}
The free energy of ATP synthesis equals the conversion of the proton partition
depth gradient $\Delta\Depth_{\ce{H+}}$ into the chemical partition depth of
ATP:
\begin{equation}
  \Delta G_{\mathrm{synthesis}}
    = RT\ln\!\left(\frac{[\mathrm{ATP}]}{[\mathrm{ADP}][\mathrm{P}_i]\,K_{\mathrm{eq}}}\right)
    = n_{\ce{H+}}\,\Delta\Depth_{\ce{H+}}\,\kB T\,\ln b,
  \label{eq:atp_synthesis}
\end{equation}
where $n_{\ce{H+}} = 3$ is the number of protons translocated per ATP
synthesised, and $\Delta\Depth_{\ce{H+}}$ is the total proton partition depth
difference across the inner mitochondrial membrane.
\end{theorem}

\begin{proof}
The proton-motive force $\Delta p$ has two components:
\begin{equation}
  \Delta p = \Delta\psi - \frac{2.303\,RT}{F}\Delta\mathrm{pH}
           \approx -180\ \mathrm{mV}\ \text{(healthy mitochondria)}.
  \label{eq:pmf}
\end{equation}
From Theorem~\ref{thm:nernst_partition}, each component corresponds to a
partition depth difference:
\begin{align}
  \Delta\psi &= \frac{\kB T}{q}\,\Delta\Depth_\psi\,\ln b, \\
  -\frac{2.303\,RT}{F}\Delta\mathrm{pH}
    &= \frac{\kB T}{q}\,\Delta\Depth_{\mathrm{pH}}\,\ln b.
\end{align}
The total proton partition depth is
$\Delta\Depth_{\ce{H+}} = \Delta\Depth_\psi + \Delta\Depth_{\mathrm{pH}}$.
For $n_{\ce{H+}} = 3$ protons per ATP:
\begin{equation}
  \Delta G = n_{\ce{H+}} \cdot q \cdot |\Delta p|
    = 3\,q\cdot\frac{\kB T}{q}\,\Delta\Depth_{\ce{H+}}\,\ln b
    = 3\,\kB T\,\Delta\Depth_{\ce{H+}}\,\ln b.
\end{equation}
Numerically, with $|\Delta p| \approx 180$\,mV:
$\Delta G = 3 \times 96485\,\mathrm{C/mol} \times 0.180\,\mathrm{V}
 \approx 52$\,kJ/mol.
This matches the observed $\Delta G_{\mathrm{ATP}} \approx 50$--$54$\,kJ/mol
under physiological conditions. Solving for
$\Delta\Depth_{\ce{H+}} \approx 7.0$ (base-$e$ units),
consistent with $|\Delta p| = 180$\,mV $= (\kB T/q)\,\Delta\Depth_{\ce{H+}}\,\ln e = 25.7\,\mathrm{mV}\times 7.0$. \qed
\end{proof}

\begin{remark}
In mammalian ATP synthase, the $c$-ring has 8 subunits, yielding a more
precise stoichiometry of $8/3 \approx 2.67$ protons per ATP when the full
$\alpha_3\beta_3$ head contributes three synthesis sites per revolution
\citep{Walker1994}. The partition converter efficiency in healthy mitochondria
is $\eta = \Delta G_{\mathrm{ATP}} / (n_{\ce{H+}} F|\Delta p|) \approx 0.85$:
85\% of the proton partition depth is captured in ATP chemical partition depth,
and 15\% is dissipated as heat.
\end{remark}

%------------------------------------------------------------------------------
\section{The \texorpdfstring{\ce{Na+}/\ce{K+}}{Na+/K+}-ATPase and Steady-State Partition Maintenance}
\label{sec:naka_atpase}
%------------------------------------------------------------------------------

The resting membrane potential is actively maintained against dissipative
leakage by the \ce{Na+}/\ce{K+}-ATPase. This pump hydrolyses one ATP to
transport three \ce{Na+} out and two \ce{K+} in per cycle, generating a net
outward current of one elementary charge per cycle (electrogenic).

In partition theory, the pump is a \emph{partition depth injector}: it
continuously invests the chemical partition depth of ATP to sustain the
partition depth differences $\Delta\Depth_{\ce{Na+}}$ and $\Delta\Depth_{\ce{K+}}$
against the dissipative tendency of leak channels.

\begin{proposition}[Steady-State Partition Depth Balance]
\label{prop:partition_balance}
At steady state, pump flux balances leak flux for each ion species:
\begin{equation}
  J_{\mathrm{pump}}^X = J_{\mathrm{leak}}^X,
  \label{eq:partition_balance}
\end{equation}
maintaining constant $\Delta\Depth_X$ and therefore a stable resting potential.
The electrogenic contribution of the pump to the resting potential is
\begin{equation}
  \Delta V_{\mathrm{pump}}
    = -\frac{I_{\mathrm{pump}}}{g_{\mathrm{total}}}
    \approx -3\ \text{to}\ -5\ \mathrm{mV},
  \label{eq:pump_hyperpolarization}
\end{equation}
accounting for the discrepancy between the GHK prediction ($\approx -65$\,mV)
and the observed resting potential ($\approx -70$\,mV).
\end{proposition}

Pump failure (ouabain blockade or ATP depletion) abolishes
Eq.~\eqref{eq:partition_balance}: $\Delta\Depth_{\ce{Na+}}$ and
$\Delta\Depth_{\ce{K+}}$ collapse over minutes, causing membrane
depolarisation, uncontrolled \ce{Ca^{2+}} influx, and excitotoxic cascades.
This cascade is analysed in Chapter~\ref{chap:charge_trajectories} as a
$\Sk$-entropy deviation pathway toward neurodegeneration.

%------------------------------------------------------------------------------
\section{Summary}
\label{sec:charge_current_summary}
%------------------------------------------------------------------------------

\begin{chapterbox}
Biological charge current is partition boundary current
(Theorem~\ref{thm:partition_boundary_current}): charge flows if and only if
the partition boundary $\partial\Part$ moves or deforms. Ion channels are
controlled partition apertures with conductance
$g = (q^2/h)\,b^{-\Depth_{\mathrm{channel}}}$
(Corollary~\ref{cor:conductance}), determined by partition topology alone.
The GHK equation emerges from partition theory with $P_X = b^{-\Depth_X}$
(Theorem~\ref{thm:ghk}); the resting potential $\approx -70$\,mV encodes the
partition depth differences $\{\Delta\Depth_X\}$ of the resting membrane.
The Nernst potential is the partition depth difference $\Delta\Depth_X$
expressed in voltage units via $V = (\kB T/zq)\,\Delta\Depth_X\,\ln b$
(Theorem~\ref{thm:nernst_partition}). Partition flux is globally conserved
(Theorem~\ref{thm:flux_conservation}), and the membrane minimises total
partition depth (Proposition~\ref{prop:membrane_minimization}). The action
potential is a cooperative partition phase transition at $\Depth_{\mathrm{crit}}$
(Theorem~\ref{thm:action_potential_transition}). ATP synthase converts the
proton partition depth gradient $\Delta\Depth_{\ce{H+}}$ into ATP chemical
partition depth at three protons per ATP with efficiency $\eta \approx 0.85$
(Theorem~\ref{thm:atp_synthesis}). The \ce{Na+}/\ce{K+}-ATPase maintains the
steady-state partition depth balance, contributing $-3$ to $-5$\,mV to the
resting potential through its electrogenic action.
\end{chapterbox}
